% ===================== ANALISIS DE RESULTADOS =====================
\section{Analisis de Resultados}
\label{sec:analisis}

% ─────────────────────────────────────────────────────────────────────────────
\subsection{Cumplimiento de Codigo (Code Compliance)}

La evaluaci\'{o}n de cumplimiento de c\'{o}digo del modelo arroj\'{o}
\textbf{CODE COMPLIANCE EVALUATION PASSED}. El esfuerzo de c\'{o}digo
m\'{a}ximo del sistema es \num{74454.4}~\si{kPa}, en el nodo 648, para el caso
2 (Alt-SUS) W+P2, frente a un admisible de \num{115142.4}~\si{kPa}; la
relaci\'{o}n de utilizaci\'{o}n m\'{a}xima es por tanto 64.7~\%. La
Tabla~\ref{tab:top5-ratio} relaciona los cinco nodos de mayor ratio del caso
cr\'{i}tico; todos cumplen con holgura el criterio Code $\leq$ Allowable.

\begin{table}[htbp]
  \centering
  \small
  \caption{Nodos de mayor ratio de esfuerzo — línea SIM-023. Criterio: Code $\leq$ Allowable.}
  \label{tab:top5-ratio}
  \begin{tabular}{crrrrrc}
    \toprule
    N & Nodo & Bending (\si{kPa}) & Code (\si{kPa}) & Allowable (\si{kPa}) & Ratio (\%) & Cumple \\
    \midrule
    1 & 648 & \num{80101.8} & \num{74454.4} & \num{115142.4} & 64.7 & Sí \\
    2 & 640 & \num{53727.3} & \num{68105.4} & \num{115142.4} & 59.1 & Sí \\
    3 & 640 & \num{53727.3} & \num{68105.4} & \num{115142.4} & 59.1 & Sí \\
    4 & 990 & \num{52166.0} & \num{66648.4} & \num{115142.4} & 57.9 & Sí \\
    5 & 990 & \num{52166.0} & \num{66648.4} & \num{115142.4} & 57.9 & Sí \\
    \bottomrule
  \end{tabular}
\end{table}


El caso que gobierna el dise\~{n}o es el sostenido alterno W+P2, es decir, la
combinaci\'{o}n de peso propio con la presi\'{o}n de operaci\'{o}n $P_2$; el
segundo caso sostenido alterno (caso 4) alcanza 58.0~\% en el nodo 990, el caso
sostenido 6 (SUS) W+P2 llega a 57.8~\% en el mismo nodo, y el caso sostenido 5
(SUS) W, sin presi\'{o}n, presenta 48.3~\% en el nodo 260. Este resultado es
caracter\'{i}stico de l\'{i}neas con tramos rectos largos entre restricciones,
donde el esfuerzo longitudinal por presi\'{o}n y peso domina sobre el rango
t\'{e}rmico: los casos de expansi\'{o}n presentan ratios moderados (4.2~\% en el
caso 7, 20.4~\% en el caso 8 y 16.2~\% en el caso 9, todos ellos en la tee del
nodo 480, donde el factor de intensificaci\'{o}n de esfuerzo amplifica la
flexi\'{o}n), frente al admisible de expansi\'{o}n aplicable en dicho nodo de
\num{236456.5}~\si{kPa}. El esfuerzo de flexi\'{o}n se concentra en la zona del
nodo 648, como evidencia el perfil de la Figura~\ref{fig:esfuerzos}, con el
m\'{a}ximo del caso cr\'{i}tico en dicho nodo (\num{80101.8}~\si{kPa} de
flexi\'{o}n). En condici\'{o}n de operaci\'{o}n (casos 1 y 3, OPE, sin
verificaci\'{o}n de c\'{o}digo) el modelo reporta un esfuerzo de flexi\'{o}n
m\'{a}ximo de \num{68550.5}~\si{kPa} en el nodo 480 (caso 3), un esfuerzo
circunferencial (\textit{hoop}) m\'{a}ximo de \num{29436.7}~\si{kPa} en el nodo
488 y una intensidad de esfuerzo m\'{a}xima de \num{80672.2}~\si{kPa} en el
nodo 480 (caso 3).

% ─────────────────────────────────────────────────────────────────────────────
\subsection{Desplazamientos y Cargas en Soportes}

Los rangos de desplazamiento son moderados: el m\'{a}ximo absoluto es
34.5~\si{mm} en la componente $DZ$ del caso 8 (EXP)
(Tabla~\ref{tab:desplaz-max}), lo que indica que la l\'{i}nea trabaja con
movimientos relativos acotados y sin riesgo de interferencias por expansi\'{o}n.
Como se advierte en la secci\'{o}n~\ref{sec:resultados}, el reporte del job no
tabula los desplazamientos absolutos de los casos sostenidos alternos 2 y 4,
limitaci\'{o}n que no afecta la verificaci\'{o}n de c\'{o}digo, pues esta se
apoya en los esfuerzos y en los rangos de expansi\'{o}n.

Las cargas en soportes (Tablas~\ref{tab:restricciones-ope} a
\ref{tab:restricciones-exp}) son moderadas: la carga significativa mayor
corresponde al soporte vertical del nodo 640 con \num{-4933}~\si{N} en el caso
sostenido y \num{-4356}~\si{N} en el caso de operaci\'{o}n; en los casos de
expansi\'{o}n las cargas sobre los soportes verticales son peque\~{n}as, con un
m\'{a}ximo de \num{2935}~\si{N} en el soporte del nodo 1186 para el caso 8. Al
no existir anclas rotacionales, las restricciones del modelo no transmiten
momentos. No obstante, la comparaci\'{o}n de las cargas en las conexiones con
los equipos (restricci\'{o}n de frontera del nodo 10 y extremos de los ramales
de los nodos 1300, 1360, 1420 y 1480) con los admisibles del fabricante queda
fuera del alcance del presente informe (secci\'{o}n~\ref{sec:alcance}).

% ─────────────────────────────────────────────────────────────────────────────
\subsection{Sintesis de Hallazgos}

La l\'{i}nea SIM-023 presenta margen frente a los criterios de
ASME B31.3-2016: el ratio m\'{a}ximo de utilizaci\'{o}n (64.7~\%) deja
disponible cerca de una tercera parte del admisible, los rangos de
desplazamiento t\'{e}rmico son moderados (m\'{a}ximo 34.5~\si{mm}) y las cargas
en soportes son moderadas. No se identifican cuellos de botella, zonas de
sobre-esfuerzo ni configuraciones de soporte cr\'{i}ticas. El aspecto pendiente
de verificaci\'{o}n externa es la aceptaci\'{o}n de las cargas en las
conexiones con los equipos (nodo 10 y extremos de los ramales de los nodos
1300, 1360, 1420 y 1480) por parte de sus fabricantes, junto con las
advertencias de modelado declaradas en la secci\'{o}n~\ref{sec:alcance}.
